针对六区域异构 AI 任务与新能源、储能、电网的耦合调度，本文构建“rolling-origin 闭环预测—语义唯一公共候选多遍搜索—来源显式储能动态规划—按情景交替搜索”的可复算框架。问题一在最终测试前以 8 个 expanding-window 原点完成 1584 次闭环评价；Macro-WAPE、Micro-WAPE 和 System Aggregate WAPE 分别为 0.8069、0.6312 和 0.2794，独立基础调度覆盖 50000 个任务。问题二构造 120000 行且语义零重复的公共候选池，五方向按相同顺序完成 6 个 pass，并以共同第 5 遍正式比较。问题三统一拆分新能源充电、网充、新能源售电和储能售电，以净电网交换定义峰值与波动，采用相邻 SOC 对动态规划真正优化净波动，并以可用 SOC 区间计算 EFC；修正后加权理想点选择平衡策略，成本、新能源率、净波动和 EFC 分别为 $-4.541740\times10^8$ 元、0.693131、22150.241 MW 和 92.401。问题四对 17 个情景统一复算能源与服务指标，并按由松到严继承碳候选；20\%、30\% 档仅标记声明搜索内未找到达标候选。两次无墙钟运行的 106 个结构化产物无语义差异，results.json 哈希均为 \texttt{c63996f2...389e}。结果只表示声明有限候选内的 scenario-feasible 解，不构成全局最优或完整 Pareto 前沿。